% =========================================================================
\section{Circuito 1: amplificador inversor}\label{sec:inversor}
% =========================================================================
Montamos el inversor de la \cref{fig:inversor}, con $R_2/R_1=10$, y lo excitamos con una sinusoide de
$1\,\mathrm{kHz}$ (canal X: $\vin$; canal Y: $\vout$).

% Figura: amplificador inversor (circuito 1) con tensiones y corrientes del análisis.
\begin{figure}[ht]
\centering
\begin{circuitikz}[circuito]
  \draw (0,0) node[op amp, ao] (A) {};
  % nudo inversor y red de realimentación
  \draw (A.-) -- ++(-0.6,0) coordinate (N);
  \draw (N) to[R, l_=$R_1$, i<^=$i_1$] ++(-2.4,0) coordinate (E);
  \draw (N) |- ++(0,1.5) coordinate (T);
  \draw (A.out) -- ++(0.7,0) coordinate (O);
  \draw (T) to[R, l=$R_2$, i>^=$i_2$] (T -| O) -- (O);
  % generador
  \draw (E) to[short, -o] ++(-0.3,0) coordinate (Ein) node[above, tens] {$\vin$};
  \draw (Ein) -- ++(-0.9,0) coordinate (G) to[sV, l_=$v_g$] ++(0,-2.0) node[ground] {};
  % entrada no inversora a tierra
  \draw (A.+) -- ++(-0.6,0) coordinate (P) -- ++(0,-0.45) node[ground] {};
  % carga y salida
  \draw (O) to[short, *-o] ++(1.0,0) coordinate (S) node[right, tens] {$\vout$};
  \draw (O) ++(0.5,0) to[R, l=$R_L$] ++(0,-2.0) node[ground] {};
  \fill (N) circle (1.6pt);
  % anotaciones del análisis
  \node[tens, anchor=east, font=\scriptsize] at ($(N)+(-0.08,0.75)$) {$\vm=0$};
  \draw[azulprof, line width=0.3pt] ($(N)+(-0.08,0.72)$) -- ($(N)+(-0.02,0.06)$);
  \node[tens, anchor=east, font=\scriptsize] at ($(P)+(-0.1,-0.25)$) {$\vp=0$};
\end{circuitikz}
\caption{Amplificador inversor (figura~2 del guion), con $R_1=1\,\mathrm{k}\Omega$,
$R_L=10\,\mathrm{k}\Omega$ y la alimentación omitida. Se indican las magnitudes del análisis: el nudo
inversor es una tierra virtual ($\vm=\vp=0$) y no absorbe corriente, de modo que $i_1=i_2$.}
\label{fig:inversor}
\end{figure}


\begin{pregunta}[Formas de onda y ganancia]
Dibujar en cada caso las formas de onda de $\Vin$ y $\Vout$, detallando las escalas de cada canal y de la base
de tiempos, y comparar la relación de amplitudes (o de tensiones pico a pico) con la $A_V$ buscada.
\end{pregunta}

\begin{respuesta}
Con el modelo ideal, la ganancia esperada, con los valores \emph{medidos} de las resistencias, es
\begin{equation}\label{eq:i-Avteo}
  A_V^{\mathrm{teo}}=-\frac{R_2}{R_1} ,
\end{equation}
y la experimental se obtiene del cociente de tensiones pico a pico, con el signo de la contrafase observada:
\begin{equation}\label{eq:i-Avexp}
  A_V^{\mathrm{exp}}=-\frac{V_{\mathrm{out},pp}}{V_{\mathrm{in},pp}},\qquad
  \varepsilon=\frac{\abs{A_V^{\mathrm{exp}}}-\abs{A_V^{\mathrm{teo}}}}{\abs{A_V^{\mathrm{teo}}}} .
\end{equation}
\begin{center}\small
\renewcommand{\arraystretch}{1.3}
\begin{tabular}{@{}cccccc@{}}
\toprule
$V_p$ nominal & Escala X, Y (V/div) & Base de tiempos & $V_{\mathrm{in},pp}$ & $V_{\mathrm{out},pp}$ & $A_V^{\mathrm{exp}}$\\
\midrule
$0{,}1\,\mathrm{V}$ & \dato & \dato & \dato & \dato & \dato\\
$0{,}2\,\mathrm{V}$ & \dato & \dato & \dato & \dato & \dato\\
$0{,}5\,\mathrm{V}$ & \dato & \dato & \dato & \dato & \dato\\
$1{,}0\,\mathrm{V}$ & \dato & \dato & \dato & \dato & \dato\\
$1{,}5\,\mathrm{V}$ & \dato & \dato & \dato & \dato & \dato\\
\bottomrule
\end{tabular}
\end{center}
\begin{center}
\hueco{0.47\textwidth}{Oscilograma $\vin$, $\vout$\\($V_p=0{,}5\,\mathrm{V}$)}\hfill
\hueco{0.47\textwidth}{Oscilograma $\vin$, $\vout$\\($V_p=1{,}5\,\mathrm{V}$)}
\end{center}
Con $A_V^{\mathrm{teo}}=$ \dato, la desviación \eqref{eq:i-Avexp} es \dato. Mientras la salida no satura, la
salida es una sinusoide en contrafase con la entrada y el cociente pico a pico reproduce \eqref{eq:i-Avteo};
las diferencias se explican por la tolerancia de las resistencias y la resolución de la lectura en pantalla.
La ganancia finita del AO no es apreciable: con $A_0\simeq2\cdot10^{5}$, la ganancia real
$\abs{A_V}=(R_2/R_1)/[1+(1+R_2/R_1)/A_0]$ difiere de la ideal en un $0{,}006$\,\%. En las amplitudes mayores
(donde $\abs{A_V}V_p>\Vsat$) el cociente deja de valer $10$ porque la salida se recorta (cuestión~2).
\end{respuesta}

\begin{pregunta}[Recorte o \emph{clipping}]
¿Qué ocurre en la forma de la salida si la amplitud de la entrada se aumenta por encima de un cierto nivel?
Medir los niveles de entrada y salida para los que el fenómeno empieza.
\end{pregunta}

\begin{respuesta}
La salida del AO no puede superar las tensiones de saturación, algo menores que la alimentación:
$-V_{\mathrm{sat}}^{-}\le\vout\le V_{\mathrm{sat}}^{+}$. La relación lineal $\vout=A_V\vin$ solo vale mientras $\abs{A_V}\,\abs{\vin}$
no alcanza esos valores; por encima, las crestas de la sinusoide de salida se recortan en tramos planos.
El recorte empieza cuando
\begin{equation}\label{eq:i-umbral}
  \abs{A_V}\,V_{p,\mathrm{umbral}}=\Vsat\quad\Longrightarrow\quad V_{p,\mathrm{umbral}}=\frac{\Vsat}{\abs{A_V}} .
\end{equation}
Medimos $V_{\mathrm{sat}}^{+}=$ \dato\ y $V_{\mathrm{sat}}^{-}=$ \dato, y el recorte comienza para una amplitud de entrada
$V_{p,\mathrm{umbral}}=$ \dato, frente a $\Vsat/\abs{A_V}=$ \dato\ según \eqref{eq:i-umbral}. Las dos saturaciones pueden
no ser iguales, porque las etapas de salida positiva y negativa del 741 no son simétricas.
\end{respuesta}

\begin{pregunta}[Característica de transferencia]
Con la amplitud de entrada máxima, trazar $\Vout(\Vin)$ en modo X--Y. Señalar para qué amplitud de entrada
(en la representación temporal) se produce el corte y relacionarlo con lo medido en la representación X--Y.
\end{pregunta}

\begin{respuesta}
En modo X--Y el osciloscopio dibuja $\vout$ frente a $\vin$ eliminando el tiempo, de modo que vemos
directamente la característica de transferencia. Esperamos tres tramos:
\begin{equation}\label{eq:i-caract}
  \vout=\begin{cases}
    +V_{\mathrm{sat}}^{+}, & \vin\le-V_{\mathrm{sat}}^{+}/\abs{A_V},\\[2pt]
    A_V\,\vin, & \text{en la zona lineal},\\[2pt]
    -V_{\mathrm{sat}}^{-}, & \vin\ge V_{\mathrm{sat}}^{-}/\abs{A_V},
  \end{cases}
\end{equation}
una recta de pendiente $A_V$ por el origen y dos tramos horizontales. La pendiente medida es \dato\ y los codos
están en $\vin=$ \dato\ y \dato.
\begin{center}
\hueco{0.6\textwidth}{Característica $\vout(\vin)$ en modo X--Y}
\end{center}
Las abscisas de los codos son exactamente las amplitudes de entrada a las que empieza el recorte en la
representación temporal \eqref{eq:i-umbral}: ambas medidas describen el mismo fenómeno, la entrada del AO en
saturación, vista en un caso frente al tiempo y en el otro frente a $\vin$.
\end{respuesta}

\begin{pregunta}[Respuesta a $100\,\mathrm{kHz}$]
¿Qué ocurre en la salida si la frecuencia de la entrada se aumenta a $100\,\mathrm{kHz}$? Representar las ondas de
entrada y salida con sus puntos más significativos. ¿Cuánto vale la ganancia?
\end{pregunta}

\begin{respuesta}
A esta frecuencia deja de valer la hipótesis $A\to\infty$, y aparecen dos limitaciones del AO real.

\textbf{Ancho de banda finito.} La ganancia en lazo abierto del 741 decae como $\abs{A}\simeq f_t/f$, con
$f_t\simeq1\,\mathrm{MHz}$. Con realimentación, el inversor se comporta como un paso bajo de primer orden de frecuencia
de corte
\begin{equation}\label{eq:i-f3dB}
  f_{3\mathrm{dB}}=\frac{f_t}{1+R_2/R_1}\simeq\frac{1\,\mathrm{MHz}}{11}\simeq91\,\mathrm{kHz},
\end{equation}
de modo que, a $f=100\,\mathrm{kHz}$,
\begin{equation}\label{eq:i-100k}
  \abs{A_V}=\frac{R_2/R_1}{\sqrt{1+(f/f_{3\mathrm{dB}})^{2}}}\simeq6{,}7,\qquad
  \Delta\varphi=-\arctg\frac{f}{f_{3\mathrm{dB}}}\simeq-48^{\circ} :
\end{equation}
la ganancia baja y la salida se retrasa, además de la inversión, unos $48^{\circ}$, es decir,
$\Delta t=\abs{\Delta\varphi}/(360^{\circ}f)\simeq1{,}3\,\mu\mathrm{s}$.

\textbf{Velocidad de subida.} La pendiente de la salida no puede superar $\mathrm{SR}\simeq0{,}5\,\mathrm{V}/\mu\mathrm{s}$. Una
sinusoide de amplitud $V_p$ y frecuencia $f$ tiene pendiente máxima $2\pi fV_p$, luego solo se reproduce sin
deformarse si
\begin{equation}\label{eq:i-slew}
  V_{p,\mathrm{out}}\le\frac{\mathrm{SR}}{2\pi f}\simeq0{,}8\,\mathrm{V}\quad(f=100\,\mathrm{kHz});
\end{equation}
por encima, la salida se convierte en una onda triangular de pendientes $\pm\mathrm{SR}$.

\textbf{Medidas.} Observamos \dato. La ganancia medida es $\abs{A_V}=$ \dato, con un retraso adicional
$\Delta t=$ \dato; los máximos de la salida valen \dato\ y aparecen en $t=$ \dato.
\begin{center}
\hueco{0.6\textwidth}{Oscilograma a $100\,\mathrm{kHz}$, con máximos, mínimos\\y sus posiciones temporales}
\end{center}
\end{respuesta}

\begin{pregunta}[Valor de las resistencias]
¿Sería razonable elegir resistencias del orden de decenas de $\Omega$? ¿Y del orden de $\mathrm{M}\Omega$?
\end{pregunta}

\begin{respuesta}
No. El modelo ideal solo fija el cociente $R_2/R_1$, pero en el laboratorio ni los valores muy pequeños ni
los muy grandes funcionan bien. Veamos qué ocurriría en la práctica en cada caso.

\textbf{Decenas de ohmios} (por ejemplo, $R_1=10\,\Omega$, $R_2=100\,\Omega$).
\begin{itemize}
  \item \emph{El AO no puede dar tanta corriente.} Para sacar $5\,\mathrm{V}$ sobre $R_2=100\,\Omega$ haría falta
        \begin{equation}\label{eq:r-Iout}
          i=\frac{\vout}{R_2}=\frac{5\,\mathrm{V}}{100\,\Omega}=50\,\mathrm{mA},
        \end{equation}
        y el 741 limita su salida a unos $25\,\mathrm{mA}$. En el osciloscopio veríamos la sinusoide de salida
        recortada mucho antes de llegar a $\pm\Vsat$, con una amplitud menor que la esperada.
  \item \emph{No se quema, pero se calienta.} El 741 tiene protección interna contra cortocircuito, así que no
        se destruye; pero mientras limita disipa del orden de $0{,}1$--$0{,}25\,\mathrm{W}$, se nota caliente y sus
        parámetros (por ejemplo, el \emph{offset}) derivan.
  \item \emph{Las resistencias sí pueden sufrir.} $R_2$ disiparía hasta $V^{2}/R=25/100=0{,}25\,\mathrm{W}$ de pico, justo
        el límite de una resistencia de $1/4\,\mathrm{W}$; con valores aún menores se quemaría.
  \item \emph{Se pierde señal en el generador.} La entrada del inversor es $R_1=10\,\Omega$, más pequeña que la
        resistencia interna del generador ($50\,\Omega$): el generador marcaría una amplitud, pero al circuito
        llegaría solo una sexta parte.
\end{itemize}

\textbf{Megaohmios} (por ejemplo, $R_1=1\,\mathrm{M}\Omega$, $R_2=10\,\mathrm{M}\Omega$).
\begin{itemize}
  \item \emph{La señal se confunde con las corrientes parásitas del AO.} Con $\vin=0{,}1\,\mathrm{V}$, la corriente de señal es
        \begin{equation}\label{eq:r-IB}
          i=\frac{\vin}{R_1}=\frac{0{,}1\,\mathrm{V}}{1\,\mathrm{M}\Omega}=100\,\mathrm{nA},
        \end{equation}
        del mismo orden que la corriente que realmente entra por las entradas del 741 ($I_B\simeq80\,\mathrm{nA}$).
        En la práctica veríamos la salida desplazada en continua: $I_BR_2\simeq0{,}8\,\mathrm{V}$ de \emph{offset}, y la
        ganancia medida ya no sería $R_2/R_1$.
  \item \emph{Más ruido e interferencias.} Los nudos de tanta resistencia actúan como antenas: la señal saldría
        con ruido y con el zumbido de $50\,\mathrm{Hz}$ de la red.
  \item \emph{Pérdida de ancho de banda.} Las capacidades parásitas del montaje (unos pocos pF) junto con
        $R_2=10\,\mathrm{M}\Omega$ forman un filtro paso bajo con corte en torno a $10\,\mathrm{kHz}$: la ganancia empezaría a
        caer cerca de la frecuencia de trabajo.
\end{itemize}

\textbf{Conclusión.} Los valores razonables están entre $1$ y $100\,\mathrm{k}\Omega$, como los de nuestro diseño
($1$ y $10\,\mathrm{k}\Omega$): corrientes de miliamperios o menos, que el AO entrega sin problema, y muy por encima
de sus corrientes de polarización.
\end{respuesta}
